\section{行政处罚}
\begin{enumerate}
  \item 具体领域的公安、环保、卫生、土地、消防等部门行政法，考试不做要求。
  \item 处罚法有 2 个法。行政处罚法（特别法）、治安管理处罚法（公安、警察、派出所字眼）。
  \item 行政机关依法对\bs{违反行政管理秩序}的公民、法人、其他组织，以\bs{减损权益 \hz 增加义务}的方式予以\bs{惩戒}。背。
  \item 特征。减损权益、增加义务（\zd{增加了新的负担}）。
        \begin{enumerate}
          \item 主体是行政机关（广义）。
          \item 以违法为前提。\al{市政府发文超 200 人以上企业缴纳 1W 环境政治费，甲对收取的 1W 不服}。
          \item 针对的是外部的公民（含外国人）、法人、其他组织。内外看身份，排除公务员。
          \item 惩戒性。目的是为了惩罚。杀鸡儆猴。
          \item \zd{vs 刑罚。1:1 折抵、未给予行政处罚的，不再给予处罚}。
        \end{enumerate}
  \item 种类。法条规定。精神、财产；资格、行为；人身、其他。
        \begin{enumerate}
          \item 精神罚。（书面）警告、（公开）通报批评。
          \item 财产罚。罚款、没收。
                \begin{kk}{没收/收缴}
                  \begin{enumerate}
                    \item 没收\zd{非法财物}。永久性的剥夺当事人的财产。
                          \begin{enumerate}
                            \item 用于违法活动且属本人的物品。\al{犯罪工具}。
                            \item 要有关联性。\al{赌徒没收宝马 \cuo}。
                            \item 属于善意第三人、受害人的不能没收。
                            \item 违法物品。原则上考试是不大会涉及到。\al{收缴非法改装车、拼接车 \dui}。
                            \item \al{决定\xhx{取缔}该公司非法采集血液行为\xhx{后}，同时没收 5 只液氮生物容器。没收 5 只液氮生物容器属于保全措施 \cuo}。
                          \end{enumerate}
                    \item \zd{不是非法财物 \yjt 享有合法所有权}。
                    \item 没收违法所得。成本 + 利润。有些成本难以切割、有些成本经非法财物转化来。
                    \item \al{没收扣押物品属于行政处罚 \dui}。
                  \end{enumerate}
                \end{kk}
          \item 资格罚。作用于许可证。暂扣许可证（\al{驾照}）、吊销许可证、降低资质等级。
          \item 行为罚。直接作用于行为。两限制两责令。
                \begin{enumerate}
                  \item 限制从业。\al{证监会对某公司负责人采取终身禁入证券市场措施}。
                  \item 限制生产经营活动。\al{原来生产 10 吨，限制生产 5 吨}。
                  \item 责令关闭。永久。
                  \item 责令停产停业。命题人认为是一段时间。
                \end{enumerate}
          \item 人身罚。行政拘留。单个 15 日。多个合并上限 20 日。只能是公安/国安来决定、执行。
          \item 其他（新花样）。
                \begin{enumerate}
                  \item 创造一些本领域特别需要的一些手段。\al{取消学籍、责令种树}。
                  \item 对外国人的责令限期离境（附加罚）。
                \end{enumerate}
        \end{enumerate}
  \item 立法。
        \begin{enumerate}
          \item 设定。倒三角形。
                \begin{enumerate}
                  \item 限制人身自由只能由法律设定。
                  \item \bs{行除人}。行政法规。除限制人身自由。
                  \item \bs{地除吊，新花样}。地方性法规。除吊销营业执照（法人是全国的）、其他种类（新花样）。
                  \item \bs{规章警告批评和罚款}。数额国务院/省人常定 \yjt 其他能自定{（\al{计算方式}）}。要能举一反三。
                \end{enumerate}
          \item \bs{具体规定}。从粗到细，下位法\zd{可以细化}上位法，但是\zd{不能抵触}上位法。\zx{扩大了、超出了}。
                \begin{enumerate}
                  \item 不能抵触上位法处罚的\zd{行为、种类、幅度}。\al{随地\xhx{吐痰}罚款 80 到 400。甲\xhx{大小便}被罚了 100}。
                  \item 细化这个区间。\al{公路法的 76 条规定擅自占用挖掘公路罚 3 万以下}。
                \end{enumerate}
          \item 补充规定/补充设定。从有到全。仅法规可以补充规定。
                \begin{enumerate}
                  \item 完整的结构 = 行为模式/条件模式 + 后果模式/责任模式。
                  \item 把缺失了的法条结构补完整。
                \end{enumerate}
        \end{enumerate}
  \item 实施。
        \begin{enumerate}
          \item 主体。具有行政处罚权的行政机关。县公安局、派出所这些。\bs{授权、委托 + 法规；集中只看省和国}。
                \begin{enumerate}
                  \item \zd{授权、委托}。+ 法规。
                        \begin{enumerate}
                          \item 得是有事业的组织，\zd{不得委托其他组织 \hz 个人}。
                          \item \zd{授权委托书。具体事项、权限、期限等内容}。
                          \item \zd{公告授权委托书}。
                        \end{enumerate}
                  \item 集中。综合执法机关（\al{城管}）。省政府和国务院决定。但人身专属公/国安。
                \end{enumerate}
          \item 管辖。级别县上，省下放；先到先得商或共。
                \begin{enumerate}
                  \item 级别。县以上政府的部门。
                  \item 地域。\al{部门规章可以对行政处罚的地域管辖进行特别规定 \dui}。
                        \begin{enumerate}
                          \item 违法发生地、结果地。
                          \item 可以下放乡镇政府、街道办 \yjt 省决定 + 要公告。
                                下放后以谁的名义实施、原部分是否可以继续实施不一定、不一定。
                                \begin{itemize}
                                  \item 授权模式 \yjt 处罚所有权发生转移。
                                  \item 委托模式 \yjt 处罚所有权保留。
                                \end{itemize}
                        \end{enumerate}
                  \item 都有管辖权。
                        \begin{enumerate}
                          \item 先到先得。
                          \item 管辖争议。协商 \yjt 共同上级指定 \hz 共同上级直接指定。
                                \al{如发生管辖权争议，由 3 个区的共同上级机关指定管辖 \cuo{以偏概全}}。
                        \end{enumerate}
                \end{enumerate}
          \item 规则。
                \begin{enumerate}
                  \item 不处罚。无责任；情节轻；初次犯、危害轻、及时改；过时效。不重载。
                        \begin{enumerate}
                          \item 无责任能力。未成年、精神疾病。
                          \item 情节轻微，主观无过错。
                          \item 初次违法 = 危害结果轻微 + 及时改正。
                          \item 过时效。\zd{治安罚 6 个月，一般 2 年，危及人身健康 5 年}。
                                \begin{itemize}
                                  \item 违法行为终了之日起算。行为的连续，不是结果的连续。\al{非法拘禁、超载车；提供假报告}。
                                  \item 只适用行政处罚！！\al{要求缴纳社会抚养\xhx{费} \cuo{行政征收}}。
                                \end{itemize}
                        \end{enumerate}
                  \item 不重罚。责令改正（强制措施），以罚代改。
                  \item \zd{不再罚}。罚款最多罚一次，就高。
                  \item 从轻 \hz 减轻。主动消除 \hz 减轻违法行为危害后果。
                  \item 从重。后果严重、教唆胁迫诱骗、打击报复、6 个月内曾受过治安管理处罚。\al{因致使未成年受到伤害，应从重处罚 \cuo}。
                  \item 电子监控。技术 + 法制部门双重审核 \yjt 行政机关二次审核 \yjt 应当向社会公布/公告。
                \end{enumerate}
          \item 程序。一总三分。
                \begin{enumerate}
                  \item 基本原则。
                        \begin{enumerate}
                          \item 2 人 + 有执法资格。
                          \item 文明执法、回避、陈述、申辩、有权要求举行听证 \bdy 应当听证/应当进行听证。
                          \item 公开有前提。社会影响大，涉密可不公开。\al{甲申请公开作出的某项处罚决定。甲申请公开的信息属于区公安分局\xhx{应当主动公开}的事项 \cuo}。
                          \item 询问、检查应当制作笔录。
                        \end{enumerate}
                  \item 简易程序 = 当场处罚。警告、一清二白、单三千。事实清。公民 200、单位 3000。拒收注明。
                  \item 普通程序。4 个环节。
                        \begin{enumerate}
                          \item 出示证件，表明身份。陈述申辩、如实陈述。
                          \item \zd{可以先行登记保存}。
                                \begin{itemize}
                                  \item \bs{证据}\zd{可能灭失 \hz 日后难以取得}。
                                  \item 经行政机关\zd{负责人批准}。\al{可以由 2 名执法人员在现场\xhx{直接}作出 \cuo}。
                                  \item 应当在\zd{7 日内}及时作出处理决定
                                \end{itemize}
                          \item 结果。负责人审查、书面决定、90 日。当场（= 当面）/7 日内送达。
                          \item Opt(专员审核)。公共利益、案情复杂。初次要过法考。
                          \item Opt(鉴定)。专员鉴定。
                          \item 法制审核。\sjt 不管啥审核。个人要求审核 = 重大利益 + 听证前置。
                        \end{enumerate}
                  \item 听证程序。应当告知听证 \bdy 应当听证/应当进行听证 \bdy 不可以要求举行听证！！！
                        \begin{enumerate}
                          \item \bs{应当告知当事人有要求听证的权利} \yjt 申请 \yjt 听证。\bs{降级、吊证、责令关、停、从业。公吊 2 千}。
                                \begin{itemize}
                                  \item \zd{降低资质等级、吊销许可证。责令关闭、责令停产停业、限制从业}。\al{暂扣、撤销许可证 \cuo}。
                                  \item 公吊 2 千。公安吊销许可证、罚款 2 千以上。\al{税务局扣押 200W 的商品 \cuo}。
                                  \item 不是有罚款就一定有听证，仔细。
                                  \item 拘留不属于法定听证情形，属于约定听证。实践中忙不过来。
                                  \item 不属于法定听证的情形，多提一嘴不算违法！！\al{某公安拟拘留甲 10 日，告知其可以申请听证 \dui}。
                                  \item \al{可以告知超市有申请听证的权利 \dui}。
                                \end{itemize}
                          \item 5207 没有 2。5 日内申请，7 日前告知 \yjt \zd{未依照法定程序召开听证}。
                          \item 申请人、第三人可以委托 1 - 2 个代理人。
                          \item \zd{主持人。非本案调查人员；直接利害关系的，当事人有权申请回避}。
                          \item 笔录。排他效力，是后续决定的唯一根据/依据。\al{应当参考 \cuo。应当根据 \dui}。
                        \end{enumerate}
                \end{enumerate}
        \end{enumerate}
  \item 治安罚。只罚个人，不罚单位。
        \begin{enumerate}
          \item 主体。县公安局、派出所。
          \item 调查。
                \begin{enumerate}
                  \item 传唤。\zx{主体是违反治安管理行为人}。\al{县公安局可以\xhx{传唤甲提供证言} \cuo}。
                        \begin{enumerate}
                          \item 负责人批准、书面、工作证 + 传唤证。
                          \item 紧急时，现场可以口头、强制传唤。告知事实理由依据。
                        \end{enumerate}
                  \item 询问。
                        \begin{enumerate}
                          \item \zd{一般 8 小时，可能拘留不超 24}。要明确提到拘留。\al{案情复杂 \cuo}。
                          \item 未成年要求监护人在场。
                          \item 证人、受害人要到其单位、住所（高效便民）。必要时可以通知来警察局。\al{如涉密}。
                        \end{enumerate}
                  \item 检查。
                        \begin{enumerate}
                          \item 两个人两个证。工作证件、县以上公安开具的检查证明。情况紧急出示工作证件即可。
                          \item 检查公民住宅，必须两证齐全。
                          \item 调查阶段，为核对证据可进行扣押（\bdy 没收）。
                        \end{enumerate}
                \end{enumerate}
          \item 决定。处罚三月七、治安二月二。
                \begin{itemize}
                  \item 期限 30 + 30。\al{县公安局对甲作出治安处罚期限为 90 日 \cuo}。
                  \item 当面/2 日内送达。要抄送家属，抄告受害人。
                \end{itemize}
          \item 执行。
                \begin{enumerate}
                  \item 拘留。
                        \begin{enumerate}
                          \item 拘留所执行。\al{看守所 \cuo}。
                          \item 暂缓执行，考试一律选不能暂缓。已申请复议 \hz 提起诉讼；主动申请；担保 \hz 保证金；公安认为暂缓执行不会发生社会危险。
                          \item 不能暂缓 \bdy 不可以提出暂缓申请。
                        \end{enumerate}
                  \item 罚款。罚执分离。可当场的 \yi 100 以下。\er 事后难以执行的。
                \end{enumerate}
          \item 从轻 \hz 减轻。14 - 18。14 不予处罚。
          \item （当事人）调解 \ydy 民诉的和解。
                \begin{enumerate}
                  \item 达成协议后。不能再进行治安处罚。
                  \item 不履行。仅此 2 法。进行行政处罚、告知提起民事诉讼。\al{申请执行调解协议书、可以裁决 \cuo}。
                  \item 调解行为。不具有强制力、不是具有处分性的具体行政行为，不可诉。
                        \al{甲不履行协议，乙可以直接向法院提起\xhx{行政诉讼} \cuo}。
                \end{enumerate}
        \end{enumerate}
\end{enumerate}
